\chapter{Conclusions}
\label{ch:conclusions}

This thesis set out to reconcile two demands that are often treated as in tension: that a teleoperated bimanual manipulator be assisted enough to be usable through a single grounded haptic interface, and that the assistance never be trusted to keep the operator safe on its own. The two demands are kept apart architecturally rather than balanced numerically --- a reactive safety filter enforces collision avoidance and joint limits as hard constraints on whatever command reaches the arms, and an alignment-based arbitration law decides, moment to moment, how much of the operator's own hand the assistance is allowed to move --- and this chapter takes stock of what evaluating that architecture, on the physical robot and at human-subject-study scale, actually shows, building on the reading of those results against expectations, the alternative explanations, and the threats to validity that Chapter~\ref{ch:discussion} already sets out.

Four design commitments carry that architecture through to evaluation. The safety filter is a single strictly convex QP, formulated natively in joint velocity rather than in the position-space updates the closest prior systems use, with a normalised CLF that lets tracking degrade gracefully under conflict while the CBF rows never do (Section~\ref{sec:clf}) --- the property the home-to-crossed hardware test of Section~\ref{sec:hwswap} is built specifically to exercise, and which it confirms directly: driven toward a target pair that is not jointly reachable, the two arms do not stall or fight the barrier, they converge on the best pose their own scheduled slack weight admits. Authority over that filter's input is arbitrated on the instantaneous alignment between the operator's twist and a belief-weighted policy, rather than on a scalar confidence or distance, so that a direct disagreement always hands full authority back to the operator; the human-subject study shows this is not merely a safe default but the mechanism participants actually prefer. The QP's own KKT shadow prices, besides certifying the barrier, are reused to schedule the per-arm tracking slack, while a separate reference limiter sitting upstream of the CLF bounds how large a decrease that row can be asked to make in a single tick, keeping the resulting tracking slack and transient inside a known range --- both mechanisms visible operating on hardware in Section~\ref{sec:hwswap}'s limiter traces and in the teleoperated trial of Section~\ref{sec:hwteleop} alike. And a grasped object passes through an explicit hybrid lifecycle --- excluded during the blind approach, relaxed-contact during closure, carried payload, and ramped reactivation on release --- each regime with its own admissible set and reset map, a lifecycle the teleoperated trial exercises through three real grasp windows, one of them a full success.

\section{Summary of findings}
\label{sec:conclfindings}

The factorial comparison isolates what actually drives objective performance: assistance, not the motion mapping, explains most of the variance in task time and path length, and within assistance the effect is carried by reference blending rather than by force guidance, which is consistently the weakest of the three settings in both mappings (Section~\ref{sec:resperf}). No significant difference separates the two blended conditions from one another, and where they interact mechanistically at all it is because a velocity mapping reads guidance force and blended twist through the same channel, an effect the position mapping does not share (Section~\ref{sec:resauthority}). Safety is the one measure on which no factor produces a detectable difference: the worst clearance recorded sits within the same few millimetres in all six conditions, consistent with the barrier setting the floor, even as the time spent actively holding that floor falls by a third under the best-assisted condition (Section~\ref{sec:ressafety}).

The subjective record complicates a purely objective reading. Participants rate the two blended conditions as more controllable and more trustworthy than force guidance alone, and rank them ahead of it almost unanimously, even though blending is precisely the channel that takes some authority away from the hand (Section~\ref{sec:ressubj}); the operators experience the arbitration as the robot doing what they mean, not as authority lost. The motion mapping shows the opposite pattern: despite separating the two objectively on smoothness, hand speed, and intent confidence, no rated item nor the composite index distinguishes the two, which is absence of evidence for a subjective difference rather than evidence of its absence. A preference for one mapping over the other, where it exists at all, does not track the measurements that are supposed to explain it.

The two hardware trials of Section~\ref{sec:hwvalidation} corroborate the same mechanisms outside simulation. The open-loop reaching test confirms that the CLF's slack yields gracefully under a genuinely infeasible joint target, holding the barrier throughout; the teleoperated trial carries a real grasp to completion under the same architecture the human-subject study evaluates, with the shared-autonomy telemetry --- belief, authority, and the grasp state machine --- behaving exactly as the simulated study's own account of those mechanisms predicts. The same trial is also where the architecture exposes the one recorded safety violation, this thesis' clearest limitation.

\section{Limitations}
\label{sec:concllimitations}

The barrier condition the QP certifies each tick is nominal: it is evaluated on the most recently available reading of the state, not a certified physical forward invariance of the sampled, disturbed, actuator-lagged closed loop. No discrete-time or robust sampled-data argument is supplied to close that gap: measurement-robust CBFs~\cite{cosner2021mrcbf} address uncertainty in the measured state but not the sampling gap itself, and a dedicated sampled-data result~\cite{breeden2021sampled} tightens the margin by a term this thesis' controller does not carry. The empirical size of the gap is what Section~\ref{sec:hwswap} and Section~\ref{sec:hwteleop} measure directly: a $-2.6$~mm margin dip recorded open-loop, and, far more severely, the teleoperated event below.

The teleoperated trial of Section~\ref{sec:hwteleop} also records the architecture's clearest limitation. During the second grasp attempt, the grasp state machine removed the gripper--target pair from the barrier for the blind insertion, by design, to permit the grasp, held it removed through the stall and a clearance-blind, fixed-duration retreat, and restored it on a timer onto an unresolved overlap. The shadow price is zero throughout the stall because the row does not exist to report anything; once re-armed, every recorded solve is feasible and satisfies its own row, yet the closed loop still fails to deliver the escape that row demands (Section~\ref{sec:cbf-scope}). The gap is therefore twofold: a discrete layer that suspends a constraint and restores it without a clearance test, and the nominal-versus-physical gap of the barrier itself once re-armed inside a violation. The current revision's clearance-gated recovery (Section~\ref{sec:grasp}) addresses the first, not the second. Chapter~\ref{ch:sota} already flags the general form of this risk --- that a formal safety guarantee at the QP level does not rule out degraded closed-loop behaviour one level of abstraction up~\cite{reis2021undesirable} --- and this trial is the first evidence, on this architecture, that the risk is not only theoretical.

Two narrower limitations follow directly from the architecture's own stated design choices. Grasp confirmation is purely kinematic, so a wedged or joint-limited retreat of the kind above is inferred from a stall rather than sensed directly (Section~\ref{sec:grasp}); and the questionnaire sample, while adequately powered for the within-subject comparisons it supports, is predominantly male and reports low prior familiarity with teleoperation devices (Section~\ref{sec:ressubj}), so the subjective preferences reported here should not be read as generalising to an experienced operator population without further evidence. The haptic loop follows standard passivity heuristics, but no passivity or stability analysis covers the complete loop, moving targets, guidance bias with cross-process latency, saturation, and slew limiting included, and the device's own damping is not characterised; stability is supported only by the absence of reported oscillation.

The adaptive slack-weight schedule of Section~\ref{sec:sched} is likewise exercised rather than isolated, since no closed-loop stability analysis of the schedule--QP--robot loop and no ablation against a fixed slack weight or a single-filter schedule are carried out, so only the two-filter design as built is evaluated here. Its sensitivity parameter $\beta$ is calibrated to one specific constraint normalisation: scaling the collision row's Jacobian and right-hand side by a common factor leaves the QP's solution unchanged but scales its dual variable inversely, so the shadow price the schedule reads is tied to that normalisation, not an invariant quantity; normalising it by the row's own gradient norm, not implemented here, would remove that dependency.

\section{Future work}
\label{sec:conclfuture}

Three directions follow from these limitations rather than from open questions in the architecture itself. The retreat--barrier interaction is the most concrete: coupling the grasp state machine's recovery motion to the same barrier gradient the QP already computes, rather than suspending that pair's barrier and restoring it on an unconditional timer, would close exactly the gap the teleoperated trial exposes, and the shadow-price burst that preceded the interpenetration is itself a plausible detector for triggering that coupled recovery before a margin is actually crossed. Tactile or force sensing at the gripper would let grasp confirmation, and the wedged-retreat case in particular, be sensed rather than inferred from a kinematic stall, complementing the purely kinematic account this thesis gives in Section~\ref{sec:grasp}. And the factorial comparison itself is the natural next replication target on physical hardware, at whatever scale a multi-session, multi-participant protocol on the real platform allows, to test whether the assistance effects Section~\ref{sec:results} reports in simulation --- and that the two hardware trials here already show operating through the same mechanisms --- persist at the same magnitude once a physical operator, not a simulated one, is driving a physical arm.

None of this reopens the trade-off the Motivation opened with. The reactive layer still never decides which object is intended or when a detour is the only way out of a dead end, and the operator's authority over any motion they oppose is never in question in the human-subject data or in the one hardware episode where it mattered most. What Section~\ref{sec:results} and Section~\ref{sec:hwvalidation} together show is a narrower and more useful thing than validation in the abstract: precision and safety are the robot's to supply, intent and judgement stay the operator's, and across a matched simulated study and two trials on physical hardware, they stay that way.

\section*{Code availability}
The software described in this thesis is available in two repositories: the robot-side stack --- safety controller, shared autonomy, grasp lifecycle, and head perception --- at \url{https://github.com/Robertorocco/triago_control}, and the operator-side stack --- device driver, teleoperation modes, and force managers --- at \url{https://github.com/Robertorocco/haption_teleoperation}.


